\documentclass[11pt]{article}
\usepackage[utf8]{inputenc}
\usepackage{amsmath, amssymb, amsthm}
\usepackage{geometry}
\geometry{letterpaper, margin=1in}
\usepackage{hyperref}
\usepackage{enumitem}

\title{Homework Chapter 3: Introduction to Determinants \\ \Large Problem Sets and Solutions}
\author{}
\date{}

\begin{document}

\maketitle
\tableofcontents
\newpage

% ---------------------------------------------------------
\section{Concept Checklist \& Question Mapping}
% ---------------------------------------------------------
This section outlines all the core concepts, problem types, and tricks found in the source material. Each concept is mapped to the corresponding problems in the Problem Set to ensure comprehensive coverage.

\begin{itemize}
    \item \textbf{C1: $2 \times 2$ and $3 \times 3$ Cofactor Expansion} (Direct computation using the recursive formula and checkerboard sign pattern). \\
    \textit{Mapped Questions:} 1, 2, 30
    
    \item \textbf{C2: Sarrus' Rule ($3 \times 3$ Diagonal Trick)} (Using the downward/upward diagonal products shortcut). \\
    \textit{Mapped Questions:} 3, 4
    
    \item \textbf{C3: Exploiting Zeros in Larger Matrices} (Scanning for rows/columns with the most zeros to collapse the sum). \\
    \textit{Mapped Questions:} 5, 6, 30
    
    \item \textbf{C4: Triangular and Diagonal Matrices} (Determinant is the product of the main diagonal). \\
    \textit{Mapped Questions:} 7, 8
    
    \item \textbf{C5: Effect of Elementary Row Operations} (Row replacement, swapping, and scaling). \\
    \textit{Mapped Questions:} 9, 10, 11
    
    \item \textbf{C6: Properties of Determinants} (Linearity pitfalls, $\det(kA) = k^n \det(A)$, $\det(AB) = \det(A)\det(B)$, $\det(A^T) = \det(A)$). \\
    \textit{Mapped Questions:} 12, 13, 14, 24, 25
    
    \item \textbf{C7: Area and Volume Applications} (Parallelograms, parallelepipeds, and transformation scaling factors). \\
    \textit{Mapped Questions:} 15, 16, 17
    
    \item \textbf{C8: Bounding and Reasonable Answers} (Finding min/max possible determinants given bounded matrix entries). \\
    \textit{Mapped Questions:} 18, 19
    
    \item \textbf{C9: Block Matrix Determinants} (Using $\det(M) = \det(A)\det(D)$ for block triangular matrices). \\
    \textit{Mapped Questions:} 20, 21
    
    \item \textbf{C10: The Vandermonde Determinant} (Polynomial interpolation matrices). \\
    \textit{Mapped Questions:} 22, 23
    
    \item \textbf{C11: Theoretical Properties and "Find the Flaw"} (True/False, definitions, common traps like the $4 \times 4$ diagonal trap). \\
    \textit{Mapped Questions:} 24, 25, 26, 27
    
    \item \textbf{C12: Abstract Real-World Modeling} (Design matrices, Jacobian matrices in Normalizing Flows). \\
    \textit{Mapped Questions:} 28, 29
\end{itemize}

\newpage

% ---------------------------------------------------------
\section{Problem Sets}
% ---------------------------------------------------------

\textbf{Directions:} Solve the following 30 problems. Show all intermediate steps, especially when applying row operations or cofactor expansions.

\subsection*{Part 1: Direct Computation and Shortcuts}

\textbf{1.} Compute the determinant of $A$ using cofactor expansion across the second row:
$$A = \begin{bmatrix} 2 & -1 & 3 \\ 0 & 4 & -2 \\ 1 & 5 & 1 \end{bmatrix}$$

\textbf{2.} Compute the determinant of $B$ using cofactor expansion down the third column:
$$B = \begin{bmatrix} 1 & 2 & 0 \\ -1 & 3 & 5 \\ 4 & -2 & 1 \end{bmatrix}$$

\textbf{3.} Compute the determinant of $C$ using Sarrus' Rule (the $3 \times 3$ diagonal trick):
$$C = \begin{bmatrix} 3 & 1 & -2 \\ 2 & 0 & 4 \\ -1 & 5 & 1 \end{bmatrix}$$

\textbf{4.} Compute the determinant of $D$ using Sarrus' Rule:
$$D = \begin{bmatrix} -2 & 4 & 1 \\ 3 & -1 & 2 \\ 0 & 2 & -3 \end{bmatrix}$$

\subsection*{Part 2: Exploiting Zeros and Matrix Structures}

\textbf{5.} Compute the determinant of the $4 \times 4$ matrix $E$ by exploiting zeros:
$$E = \begin{bmatrix} 2 & 0 & 0 & 3 \\ -1 & 5 & 0 & 4 \\ 3 & -2 & 0 & -1 \\ 4 & 1 & 2 & 2 \end{bmatrix}$$

\textbf{6.} Compute the determinant of the $5 \times 5$ matrix $F$. Choose your expansions strategically:
$$F = \begin{bmatrix} 1 & 2 & -1 & 0 & 3 \\ 0 & 4 & 0 & 0 & 0 \\ 5 & -3 & 2 & 0 & 1 \\ 0 & 1 & -2 & 3 & 4 \\ 2 & 0 & 1 & 0 & -1 \end{bmatrix}$$

\textbf{7.} Compute the determinant of the triangular matrix $G$:
$$G = \begin{bmatrix} 5 & 2 & -3 & 1 \\ 0 & -2 & 4 & 7 \\ 0 & 0 & 3 & -1 \\ 0 & 0 & 0 & -1 \end{bmatrix}$$

\textbf{8.} Compute the determinant of the diagonal matrix $H = \text{diag}(2, -1, 4, 0.5, -3)$.

\subsection*{Part 3: Row Operations and Properties}

\textbf{9.} Let $A$ be a $3 \times 3$ matrix with $\det(A) = 7$. Matrix $B$ is obtained by swapping row 1 and row 3 of $A$, and then multiplying row 2 by 4. What is $\det(B)$?

\textbf{10.} Let $C$ be a $4 \times 4$ matrix with $\det(C) = -4$. Matrix $D$ is obtained by adding 3 times row 1 to row 4, then swapping row 2 and row 3, and finally swapping row 1 and row 4. What is $\det(D)$?

\textbf{11.} Use elementary row operations to simplify matrix $J$ to an upper triangular matrix, then compute its determinant:
$$J = \begin{bmatrix} 1 & 2 & 3 & 4 \\ 1 & 3 & 5 & 7 \\ 1 & 2 & 4 & 6 \\ 1 & 2 & 3 & 5 \end{bmatrix}$$

\textbf{12.} Let $A$ be a $3 \times 3$ matrix with $\det(A) = 5$. Find $\det(2A)$.

\textbf{13.} Let $A$ and $B$ be $4 \times 4$ matrices with $\det(A) = -2$ and $\det(B) = 3$. Find $\det(-A B^T)$.

\textbf{14.} Let $A$ be an invertible $3 \times 3$ matrix with $\det(A) = 4$. Find $\det(A^{-1} A^2)$.

\subsection*{Part 4: Area, Volume, and Transformations}

\textbf{15.} Find the area of the parallelogram in $\mathbb{R}^2$ with vertices at $(0,0), (4,1), (2,5),$ and $(6,6)$.

\textbf{16.} Find the volume of the parallelepiped in $\mathbb{R}^3$ determined by the vectors $\mathbf{u} = \begin{bmatrix} 1 \\ 0 \\ 2 \end{bmatrix}$, $\mathbf{v} = \begin{bmatrix} 0 \\ 3 \\ -1 \end{bmatrix}$, and $\mathbf{w} = \begin{bmatrix} 4 \\ 1 \\ 1 \end{bmatrix}$.

\textbf{17.} A linear transformation $T(\mathbf{x}) = A\mathbf{x}$ maps a triangle of area 5 to a region of area 35. If $A = \begin{bmatrix} 2 & k \\ -1 & 4 \end{bmatrix}$, find all possible real values of $k$.

\subsection*{Part 5: Bounding and Advanced Matrix Types}

\textbf{18.} Let $A$ be a $2 \times 2$ matrix where every entry is a real number between $-5$ and $5$ inclusive. What is the maximum possible value for $\det(A)$?

\textbf{19.} Let $B$ be a $3 \times 3$ matrix where every entry is either $0$ or $1$. What is the maximum possible value for $\det(B)$? Provide an example matrix that achieves this determinant.

\textbf{20.} Compute the determinant of the block upper-triangular matrix $M$:
$$M = \begin{bmatrix} 2 & 1 & 5 & 6 \\ -1 & 3 & 7 & 8 \\ 0 & 0 & 4 & -2 \\ 0 & 0 & 1 & 1 \end{bmatrix}$$

\textbf{21.} Compute the determinant of the block lower-triangular matrix $N$:
$$N = \begin{bmatrix} 3 & 2 & 0 & 0 & 0 \\ 1 & 4 & 0 & 0 & 0 \\ 5 & 6 & 1 & -1 & 2 \\ -2 & 3 & 0 & 2 & 3 \\ 1 & 1 & 0 & 0 & -1 \end{bmatrix}$$

\textbf{22.} Compute the determinant of the Vandermonde matrix $V_1$:
$$V_1 = \begin{bmatrix} 1 & 2 & 4 \\ 1 & 3 & 9 \\ 1 & 5 & 25 \end{bmatrix}$$

\textbf{23.} Compute the determinant of the Vandermonde matrix $V_2$:
$$V_2 = \begin{bmatrix} 1 & -1 & 1 \\ 1 & 2 & 4 \\ 1 & 4 & 16 \end{bmatrix}$$

\subsection*{Part 6: Diagnostic Testing ("Find the Flaw") and Theory}

\textbf{24.} \textbf{Find the Flaw:} A student is asked to compute $\det(A+B)$ where $A = \begin{bmatrix} 1 & 0 \\ 0 & 1 \end{bmatrix}$ and $B = \begin{bmatrix} -1 & 0 \\ 0 & -1 \end{bmatrix}$. The student writes: "$\det(A+B) = \det(A) + \det(B) = 1 + 1 = 2$." Identify the error and provide the correct solution.

\textbf{25.} \textbf{Find the Flaw:} A student is asked to compute $\det(3A)$ where $A$ is a $2 \times 2$ matrix with $\det(A) = 4$. The student writes: "Since the determinant is a linear operator, $\det(3A) = 3\det(A) = 3(4) = 12$." Identify the error and provide the correct solution.

\textbf{26.} \textbf{True or False:} If $A$ is a $4 \times 4$ matrix, Sarrus' Rule can be used to compute its determinant by copying the first three columns to the right of the matrix and summing the diagonal products. Explain your reasoning.

\textbf{27.} \textbf{True or False:} The $(i,j)$-cofactor of a matrix $A$ is defined as the submatrix obtained by deleting the $i$-th row and $j$-th column. Explain your reasoning.

\subsection*{Part 7: Real-World Modeling and Abstract Applications}

\textbf{28.} In design of experiments, we often want to maximize $\det(D^T D)$. Given the design matrix $D = \begin{bmatrix} 1 & 1 & 1 \\ 1 & -1 & 1 \\ -1 & 1 & 1 \end{bmatrix}$, compute $\det(D^T D)$.

\textbf{29.} In Generative AI (Normalizing Flows), the Jacobian matrix $J$ of a neural network transformation is strictly designed to be lower triangular to avoid $O(n^3)$ computation time. Given the Jacobian matrix below, compute the exact log-likelihood adjustment, which is $\ln(\det(J))$:
$$J = \begin{bmatrix} e^{x_1} & 0 & 0 \\ \sin(x_1) & e^{x_2} & 0 \\ x_1 x_2 & \cos(x_2) & e^{x_3} \end{bmatrix}$$

\textbf{30.} Compute the determinant of the following $4 \times 4$ matrix $K$ in terms of variables $x$ and $y$. (Hint: Use cofactor expansion down the first column).
$$K = \begin{bmatrix} x & y & 0 & 0 \\ 0 & x & y & 0 \\ 0 & 0 & x & y \\ y & 0 & 0 & x \end{bmatrix}$$

\newpage

% ---------------------------------------------------------
\section{Detailed Solutions}
% ---------------------------------------------------------

\textbf{Solution 1:}
We expand across the second row ($i=2$). The signs for the second row follow the pattern $-, +, -$.
$$\det A = -0 \cdot \det \begin{bmatrix} -1 & 3 \\ 5 & 1 \end{bmatrix} + 4 \cdot \det \begin{bmatrix} 2 & 3 \\ 1 & 1 \end{bmatrix} - (-2) \cdot \det \begin{bmatrix} 2 & -1 \\ 1 & 5 \end{bmatrix}$$
$$= 0 + 4((2)(1) - (3)(1)) + 2((2)(5) - (-1)(1))$$
$$= 4(2 - 3) + 2(10 + 1) = 4(-1) + 2(11) = -4 + 22 = 18$$
\textbf{Final Answer:} $18$

\vspace{0.5cm}
\textbf{Solution 2:}
We expand down the third column ($j=3$). The signs for the third column are $+, -, +$.
$$\det B = 0 \cdot \det \begin{bmatrix} -1 & 3 \\ 4 & -2 \end{bmatrix} - 5 \cdot \det \begin{bmatrix} 1 & 2 \\ 4 & -2 \end{bmatrix} + 1 \cdot \det \begin{bmatrix} 1 & 2 \\ -1 & 3 \end{bmatrix}$$
$$= 0 - 5((1)(-2) - (2)(4)) + 1((1)(3) - (2)(-1))$$
$$= -5(-2 - 8) + 1(3 + 2) = -5(-10) + 5 = 50 + 5 = 55$$
\textbf{Final Answer:} $55$

\vspace{0.5cm}
\textbf{Solution 3:}
Using Sarrus' Rule, we copy the first two columns:
$$
\begin{matrix}
3 & 1 & -2 & | & 3 & 1 \\
2 & 0 & 4 & | & 2 & 0 \\
-1 & 5 & 1 & | & -1 & 5
\end{matrix}
$$
Sum of downward diagonals ($D$):
$D = (3)(0)(1) + (1)(4)(-1) + (-2)(2)(5) = 0 - 4 - 20 = -24$ \\
Sum of upward diagonals ($U$):
$U = (-1)(0)(-2) + (5)(4)(3) + (1)(2)(1) = 0 + 60 + 2 = 62$ \\
$\det C = D - U = -24 - 62 = -86$ \\
\textbf{Final Answer:} $-86$

\vspace{0.5cm}
\textbf{Solution 4:}
Using Sarrus' Rule:
$$
\begin{matrix}
-2 & 4 & 1 & | & -2 & 4 \\
3 & -1 & 2 & | & 3 & -1 \\
0 & 2 & -3 & | & 0 & 2
\end{matrix}
$$
Downward ($D$): $(-2)(-1)(-3) + (4)(2)(0) + (1)(3)(2) = -6 + 0 + 6 = 0$ \\
Upward ($U$): $(0)(-1)(1) + (2)(2)(-2) + (-3)(3)(4) = 0 - 8 - 36 = -44$ \\
$\det D = D - U = 0 - (-44) = 44$ \\
\textbf{Final Answer:} $44$

\vspace{0.5cm}
\textbf{Solution 5:}
Column 3 has three zeros. We expand down Column 3. The only non-zero entry is $2$, located at $(4,3)$. Its sign is $(-1)^{4+3} = -1$.
$$\det E = -2 \cdot \det \begin{bmatrix} 2 & 0 & 3 \\ -1 & 5 & 4 \\ 3 & -2 & -1 \end{bmatrix}$$
Now compute the $3 \times 3$ determinant (expanding across row 1):
$$2((5)(-1) - (4)(-2)) - 0 + 3((-1)(-2) - (5)(3)) = 2(-5 + 8) + 3(2 - 15) = 2(3) + 3(-13) = 6 - 39 = -33$$
Multiply by the factored constant: $\det E = -2(-33) = 66$. \\
\textbf{Final Answer:} $66$

\vspace{0.5cm}
\textbf{Solution 6:}
Row 2 has four zeros. Expand across Row 2. The entry $4$ is at $(2,2)$, sign is $(-1)^{2+2} = +1$.
$$\det F = 4 \cdot \det \begin{bmatrix} 1 & -1 & 0 & 3 \\ 5 & 2 & 0 & 1 \\ 0 & -2 & 3 & 4 \\ 2 & 1 & 0 & -1 \end{bmatrix}$$
In this $4 \times 4$ matrix, Column 3 has three zeros. Expand down Column 3. The entry $3$ is at $(3,3)$, sign is $(-1)^{3+3} = +1$.
$$\det F = 4 \cdot 3 \cdot \det \begin{bmatrix} 1 & -1 & 3 \\ 5 & 2 & 1 \\ 2 & 1 & -1 \end{bmatrix}$$
Compute the $3 \times 3$ determinant (using Sarrus or cofactor):
$1(-2 - 1) - (-1)(-5 - 2) + 3(5 - 4) = 1(-3) + 1(-7) + 3(1) = -3 - 7 + 3 = -7$.
Total determinant: $4 \cdot 3 \cdot (-7) = -84$. \\
\textbf{Final Answer:} $-84$

\vspace{0.5cm}
\textbf{Solution 7:}
Matrix $G$ is an upper triangular matrix. The determinant is the product of its main diagonal entries.
$\det G = (5)(-2)(3)(-1) = 30$. \\
\textbf{Final Answer:} $30$

\vspace{0.5cm}
\textbf{Solution 8:}
Matrix $H$ is a diagonal matrix (a special case of triangular).
$\det H = (2)(-1)(4)(0.5)(-3) = 12$. \\
\textbf{Final Answer:} $12$

\vspace{0.5cm}
\textbf{Solution 9:}
Swapping two rows multiplies the determinant by $-1$. Multiplying a row by a scalar $k$ multiplies the determinant by $k$.
$\det B = \det(A) \cdot (-1) \cdot (4) = 7 \cdot (-4) = -28$. \\
\textbf{Final Answer:} $-28$

\vspace{0.5cm}
\textbf{Solution 10:}
Adding a multiple of one row to another (row replacement) does not change the determinant (factor of $1$). Swapping rows multiplies by $-1$. There are two row swaps.
$\det D = \det(C) \cdot (1) \cdot (-1) \cdot (-1) = -4 \cdot 1 = -4$. \\
\textbf{Final Answer:} $-4$

\vspace{0.5cm}
\textbf{Solution 11:}
Apply row replacements to create zeros below the pivot in column 1:
$R_2 \leftarrow R_2 - R_1$, $R_3 \leftarrow R_3 - R_1$, $R_4 \leftarrow R_4 - R_1$.
$$J \sim \begin{bmatrix} 1 & 2 & 3 & 4 \\ 0 & 1 & 2 & 3 \\ 0 & 0 & 1 & 2 \\ 0 & 0 & 0 & 1 \end{bmatrix}$$
Row replacements do not change the determinant. The new matrix is upper triangular.
$\det J = (1)(1)(1)(1) = 1$. \\
\textbf{Final Answer:} $1$

\vspace{0.5cm}
\textbf{Solution 12:}
Scaling an $n \times n$ matrix by a constant $k$ scales the determinant by $k^n$. Here $n=3$ and $k=2$.
$\det(2A) = 2^3 \det(A) = 8(5) = 40$. \\
\textbf{Final Answer:} $40$

\vspace{0.5cm}
\textbf{Solution 13:}
Using properties of determinants: $\det(kA) = k^n \det(A)$, $\det(AB) = \det(A)\det(B)$, and $\det(B^T) = \det(B)$.
$\det(-A B^T) = (-1)^4 \det(A) \det(B^T) = 1 \cdot (-2) \cdot (3) = -6$. \\
\textbf{Final Answer:} $-6$

\vspace{0.5cm}
\textbf{Solution 14:}
$\det(A^{-1} A^2) = \det(A^{-1}) \det(A^2) = \frac{1}{\det(A)} \cdot (\det(A))^2 = \det(A) = 4$. \\
\textbf{Final Answer:} $4$

\vspace{0.5cm}
\textbf{Solution 15:}
The parallelogram is spanned by the vectors $\mathbf{u} = \begin{bmatrix} 4 \\ 1 \end{bmatrix}$ and $\mathbf{v} = \begin{bmatrix} 2 \\ 5 \end{bmatrix}$.
The area is the absolute value of the determinant of the matrix formed by these vectors.
$A = \begin{bmatrix} 4 & 2 \\ 1 & 5 \end{bmatrix}$. $\det A = (4)(5) - (2)(1) = 20 - 2 = 18$.
Area = $|18| = 18$. \\
\textbf{Final Answer:} $18$

\vspace{0.5cm}
\textbf{Solution 16:}
The volume of the parallelepiped is the absolute value of the determinant of the $3 \times 3$ matrix formed by the vectors.
$V = \begin{bmatrix} 1 & 0 & 4 \\ 0 & 3 & 1 \\ 2 & -1 & 1 \end{bmatrix}$.
Expand across row 1:
$\det V = 1((3)(1) - (1)(-1)) - 0 + 4((0)(-1) - (3)(2)) = 1(3 + 1) + 4(0 - 6) = 4 - 24 = -20$.
Volume = $|-20| = 20$. \\
\textbf{Final Answer:} $20$

\vspace{0.5cm}
\textbf{Solution 17:}
The scaling factor of the area is given by $|\det A|$.
$|\det A| = \frac{\text{New Area}}{\text{Old Area}} = \frac{35}{5} = 7$.
Compute $\det A = (2)(4) - (k)(-1) = 8 + k$.
We have $|8 + k| = 7$. This gives two equations:
$8 + k = 7 \implies k = -1$
$8 + k = -7 \implies k = -15$ \\
\textbf{Final Answer:} $k = -1, -15$

\vspace{0.5cm}
\textbf{Solution 18:}
Let $A = \begin{bmatrix} a & b \\ c & d \end{bmatrix}$. $\det A = ad - bc$.
To maximize $\det A$, we must maximize $ad$ and minimize $bc$.
Since $-5 \le a,b,c,d \le 5$, the maximum value for $ad$ is $(5)(5) = 25$ or $(-5)(-5) = 25$.
The minimum value for $bc$ is $(5)(-5) = -25$.
Max determinant = $25 - (-25) = 50$. \\
\textbf{Final Answer:} $50$

\vspace{0.5cm}
\textbf{Solution 19:}
For a $3 \times 3$ matrix with entries $0$ or $1$, the maximum determinant is $2$.
An example matrix is $B = \begin{bmatrix} 0 & 1 & 1 \\ 1 & 0 & 1 \\ 1 & 1 & 0 \end{bmatrix}$.
Using Sarrus' Rule: $D = 0 + 1 + 1 = 2$, $U = 0 + 0 + 0 = 0$. $\det B = 2 - 0 = 2$. \\
\textbf{Final Answer:} $2$ (Example provided above)

\vspace{0.5cm}
\textbf{Solution 20:}
Matrix $M$ is a block upper-triangular matrix: $M = \begin{bmatrix} A & B \\ 0 & D \end{bmatrix}$, where $A = \begin{bmatrix} 2 & 1 \\ -1 & 3 \end{bmatrix}$ and $D = \begin{bmatrix} 4 & -2 \\ 1 & 1 \end{bmatrix}$.
$\det M = \det(A) \cdot \det(D)$.
$\det A = (2)(3) - (1)(-1) = 6 + 1 = 7$.
$\det D = (4)(1) - (-2)(1) = 4 + 2 = 6$.
$\det M = 7 \cdot 6 = 42$. \\
\textbf{Final Answer:} $42$

\vspace{0.5cm}
\textbf{Solution 21:}
Matrix $N$ is a block lower-triangular matrix: $N = \begin{bmatrix} A & 0 \\ C & D \end{bmatrix}$, where $A = \begin{bmatrix} 3 & 2 \\ 1 & 4 \end{bmatrix}$ and $D = \begin{bmatrix} 1 & -1 & 2 \\ 0 & 2 & 3 \\ 0 & 0 & -1 \end{bmatrix}$.
$\det N = \det(A) \cdot \det(D)$.
$\det A = (3)(4) - (2)(1) = 12 - 2 = 10$.
Matrix $D$ is upper triangular, so $\det D = (1)(2)(-1) = -2$.
$\det N = 10 \cdot (-2) = -20$. \\
\textbf{Final Answer:} $-20$

\vspace{0.5cm}
\textbf{Solution 22:}
$V_1$ is a $3 \times 3$ Vandermonde matrix with $x_1 = 2, x_2 = 3, x_3 = 5$.
The formula is $\det V = (x_2 - x_1)(x_3 - x_1)(x_3 - x_2)$.
$\det V_1 = (3 - 2)(5 - 2)(5 - 3) = (1)(3)(2) = 6$. \\
\textbf{Final Answer:} $6$

\vspace{0.5cm}
\textbf{Solution 23:}
$V_2$ is a $3 \times 3$ Vandermonde matrix with $x_1 = -1, x_2 = 2, x_3 = 4$.
$\det V_2 = (2 - (-1))(4 - (-1))(4 - 2) = (3)(5)(2) = 30$. \\
\textbf{Final Answer:} $30$

\vspace{0.5cm}
\textbf{Solution 24:}
\textbf{Error:} The student assumed the determinant is linear over matrix addition, i.e., $\det(A+B) = \det(A) + \det(B)$, which is generally false.
\textbf{Correct Solution:} First, compute the matrix sum: $A + B = \begin{bmatrix} 1-1 & 0 \\ 0 & 1-1 \end{bmatrix} = \begin{bmatrix} 0 & 0 \\ 0 & 0 \end{bmatrix}$.
Then compute the determinant: $\det(A+B) = \det(0) = 0$.

\vspace{0.5cm}
\textbf{Solution 25:}
\textbf{Error:} The student assumed $\det(kA) = k \det(A)$. However, scaling an $n \times n$ matrix scales every row, so the constant is pulled out $n$ times.
\textbf{Correct Solution:} For a $2 \times 2$ matrix, $\det(3A) = 3^2 \det(A) = 9(4) = 36$.

\vspace{0.5cm}
\textbf{Solution 26:}
\textbf{False.} Sarrus' Rule (the diagonal trick) ONLY works for $3 \times 3$ matrices. Applying it to a $4 \times 4$ matrix is a fatal error that yields mathematically invalid results. Cofactor expansion or row reduction must be used for $4 \times 4$ matrices.

\vspace{0.5cm}
\textbf{Solution 27:}
\textbf{False.} The matrix obtained by deleting the $i$-th row and $j$-th column is called a \textit{submatrix} ($A_{ij}$). The $(i,j)$-cofactor is a \textit{scalar number} defined as $C_{ij} = (-1)^{i+j} \det(A_{ij})$. It is not a matrix itself.

\vspace{0.5cm}
\textbf{Solution 28:}
Using the property $\det(D^T D) = \det(D^T)\det(D) = (\det D)^2$.
First, compute $\det D$:
$\det D = 1((-1)(1) - (1)(1)) - 1((1)(1) - (1)(-1)) + 1((1)(1) - (-1)(-1))$
$= 1(-2) - 1(2) + 1(0) = -2 - 2 = -4$.
Then, $\det(D^T D) = (-4)^2 = 16$. \\
\textbf{Final Answer:} $16$

\vspace{0.5cm}
\textbf{Solution 29:}
Because $J$ is a lower triangular matrix, its determinant is simply the product of its main diagonal entries (Theorem 2).
$\det J = (e^{x_1})(e^{x_2})(e^{x_3}) = e^{x_1 + x_2 + x_3}$.
The log-likelihood adjustment is $\ln(\det J) = \ln(e^{x_1 + x_2 + x_3}) = x_1 + x_2 + x_3$. \\
\textbf{Final Answer:} $x_1 + x_2 + x_3$

\vspace{0.5cm}
\textbf{Solution 30:}
We expand down the first column. The non-zero entries are $x$ at $(1,1)$ and $y$ at $(4,1)$.
The sign for $(1,1)$ is $+1$. The sign for $(4,1)$ is $(-1)^{4+1} = -1$.
$$\det K = x \cdot \det \begin{bmatrix} x & y & 0 \\ 0 & x & y \\ 0 & 0 & x \end{bmatrix} - y \cdot \det \begin{bmatrix} y & 0 & 0 \\ x & y & 0 \\ 0 & x & y \end{bmatrix}$$
The first $3 \times 3$ matrix is upper triangular, so its determinant is $x \cdot x \cdot x = x^3$.
The second $3 \times 3$ matrix is lower triangular, so its determinant is $y \cdot y \cdot y = y^3$.
Substituting these back:
$\det K = x(x^3) - y(y^3) = x^4 - y^4$. \\
\textbf{Final Answer:} $x^4 - y^4$

\end{document}